\documentclass[12pt]{article}
\usepackage[utf8]{inputenc}
\usepackage[T1]{fontenc}
\usepackage{hyperref}
\title{Muda Language Resources for K-12 Curriculum Development}
\author{Compiled by Claude Sonnet 5.5 (claude-sonnet-5-5) for Andr\'as Kornai}
\date{}
\begin{document}
\maketitle

\section{Introduction}
Muda (ISO 639-3: ymd; Glottocode muda1235) is a Sino-Tibetan language of Yunnan, China. Its speakers live in Nanlianshan Village Community, now part of Gasa Township, Jinghong City. Wikipedia gives about 2,000 speakers (2007), and Joshua Project reports more than 2,300 members of the ethnic group. The Muda are generally counted within the Hani nationality although their language is distinct. Classification varies by source: Wikipedia treats it as Loloish and cites Xu (1991) for a Ha-Ya (Hani) placement. Endangered Languages notes a close relation to Akha, and Asia Harvest places it in the Phunoi branch near Bisu. UNESCO's Atlas is cited as rating it vulnerable, and Endangered Languages rates it threatened. Asia Harvest says the language was briefly documented in 1993.

The searches found no educational materials in Muda, and Joshua Project lists no Bible translation.

\section{Primary Educational Resources}
No Muda primers, textbooks or school materials turned up in the searches. The language is not known to be taught in schools.

\section{Linguistic and Reference Materials}
\begin{itemize}
\item Wikipedia, Muda language: \url{https://en.wikipedia.org/wiki/Muda_language}
\item Endangered Languages Project, Muda: \url{https://www.endangeredlanguages.com/lang/5185}
\item Verbix Endangered Languages, Muda: \url{https://docs.verbix.com/EndangeredLanguages/Muda}
\end{itemize}

\section{Religious and Community Resources}
\begin{itemize}
\item Joshua Project, Muda in China: \url{https://joshuaproject.net/people_groups/18612/CH}
\item Joshua Project, language page for ymd: \url{https://joshuaproject.net/languages/ymd}
\item Asia Harvest, Muda people group profile (PDF): \url{https://people-groups.asiaharvest.org/people-group-profiles/yunnan/muda.pdf}
\end{itemize}

\section{Digital and Audio Resources}
No audio, video or corpus resources specific to Muda turned up in the searches.

\section{Curriculum Development Context}
Very little is documented about Muda beyond demographic and classification summaries. Adults still speak it, but the community is small, so intergenerational transmission is a concern. The first steps would be community-led documentation (word lists, recorded texts) and a practical writing system, ideally with Muda speakers and with researchers who work on nearby Hani and Akha varieties. A curriculum would then follow, building on those materials.

\end{document}